% ============================================================================
%  Federated and meta-learned blood-glucose forecasting across sites
%  PLOS Digital Health submission (PLOS template v3.8 Apr 2026)
%
%  STATUS: Full coherent draft. Abstract, Author summary, Introduction,
%  Materials and methods, Results (incl. finetuning), and Discussion (incl.
%  external-comparison subsection) written. Under reviewer-style revision.
%
%  NOTE: this uses the PLOS class options and the `cite` package (numeric
%  Vancouver citations). The bibliography style is `plos2015` (the genuine PLOS
%  Vancouver BibTeX style; the current template ships an identical file renamed
%  plos2025.bst). plos2015.bst is committed alongside this .tex so it compiles
%  as-is on Overleaf. The bib database is references.bib.
% ============================================================================
\documentclass[10pt,letterpaper]{article}
\usepackage[top=0.85in,left=2.75in,footskip=0.75in]{geometry}

\usepackage{amsmath,amssymb}
\usepackage{changepage}
\usepackage{textcomp,marvosym}
\usepackage{cite}
\usepackage{nameref,hyperref}
\usepackage[right]{lineno}
\usepackage[nopatch=eqnum]{microtype}
\DisableLigatures[f]{encoding = *, family = * }
\usepackage[table]{xcolor}
\usepackage{array}

\newcolumntype{+}{!{\vrule width 2pt}}
\newlength\savedwidth
\newcommand\thickcline[1]{%
  \noalign{\global\savedwidth\arrayrulewidth\global\arrayrulewidth 2pt}%
  \cline{#1}%
  \noalign{\vskip\arrayrulewidth}%
  \noalign{\global\arrayrulewidth\savedwidth}%
}
\newcommand\thickhline{\noalign{\global\savedwidth\arrayrulewidth\global\arrayrulewidth 2pt}%
\hline
\noalign{\global\arrayrulewidth\savedwidth}}

% Text layout
\raggedright
\setlength{\parindent}{0.5cm}
\textwidth 5.25in
\textheight 8.75in

\usepackage[aboveskip=1pt,labelfont=bf,labelsep=period,justification=raggedright,singlelinecheck=off]{caption}
\renewcommand{\figurename}{Fig}

\bibliographystyle{plos2015}

\makeatletter
\renewcommand{\@biblabel}[1]{\quad#1.}
\makeatother

% Header and Footer with logo
\usepackage{lastpage,fancyhdr,graphicx}
\usepackage{epstopdf}
\pagestyle{fancy}
\fancyhf{}
\rfoot{\thepage/\pageref{LastPage}}
\renewcommand{\headrulewidth}{0pt}
\renewcommand{\footrule}{\hrule height 2pt \vspace{2mm}}
\fancyheadoffset[L]{2.25in}
\fancyfootoffset[L]{2.25in}
\lfoot{\today}

\begin{document}
\vspace*{0.2in}

% Title must be 250 characters or less. Sentence case.
\begin{flushleft}
{\Large
\textbf{Federated and meta-learned blood-glucose forecasting across sites: a domain-generalisation benchmark and a cross-site deployment prototype}
}
\newline
\\
Name1 Surname\textsuperscript{1*}
\\
\bigskip
\textbf{1} Affiliation Dept/Program/Center, Institution Name, City, State, Country
\\
\bigskip

* correspondingauthor@institute.edu

\end{flushleft}

% Please keep the abstract below 300 words
\section*{Abstract}
Algorithms that predict blood glucose 30 minutes ahead are increasingly used in
type~1 diabetes (T1D) to warn of excursions and support insulin decisions, but a
model trained on one cohort often transfers poorly to new patients, devices, and
sites --- and the diverse, multi-site data that would make it robust are kept
siloed by privacy regulation. Federated learning (FL) lets institutions train a
shared model without moving raw records. Yet for glucose forecasting the leading
federated strategies have rarely been compared head-to-head, evaluations seldom
separate accuracy on the training cohorts (held-in) from accuracy on entirely
unseen cohorts (out-of-distribution, OOD), and almost all such work is simulated
rather than run across real machines. We address these gaps. Using a common
encoder--decoder transformer and four real T1D cohorts as federated clients ---
with three never-seen cohorts (ReplaceBG, BrisT1D, Flair) as OOD test sets --- we
benchmark
single-cohort local training against federated averaging (FedAvg), FedProx,
meta-learning for domain generalisation (MLDG), and two personalised-FL methods,
under a strict by-patient split and a five-seed protocol, reporting held-in and
OOD error separately at the 30-minute horizon. Federation matched local training
in-domain (held-in average $\approx$20 mg/dL RMSE) and transferred more reliably
out of it: on all three unseen cohorts the global models beat the mean
single-cohort model, and a model trained on the smallest
cohort alone collapsed on every one of them (e.g.\ $28.5$ vs $\approx$21.3 mg/dL
on ReplaceBG). The gains over local training are small relative to CGM sensor
error, so the practical value is reduced deployment risk rather than a large
accuracy jump; MLDG was numerically best in every comparison but by margins
within seed-to-seed variation (not significant at five seeds). Finetuning the federated model on a new site
did slightly better than training there from scratch. We also built a working
lightweight FL system --- a stateless server plus per-cohort clients
communicating over HTTP with only the Python standard library --- and ran it
end-to-end across two physically separate machines over the public internet,
rather than simulating federation as data shards on one host as almost all prior
glucose work does. The contribution is thus a clean, by-patient, held-in-versus-
OOD benchmark of federated strategies across four real T1D cohorts, together with
a runnable cross-machine federation: a privacy-preserving way to reduce the risk
of catastrophic transfer failure, using tooling light enough for a small clinical
consortium.

% Author Summary: 150-200 words, first person.
\section*{Author summary}
People with type~1 diabetes rely increasingly on algorithms that predict where
their glucose will be in the next half hour. These predictors work well for the
patients they were trained on but can fail when applied to a different clinic,
device, or population --- and the varied data that would make them robust are
scattered across institutions that, for privacy reasons, cannot pool raw
records. Federated learning lets those institutions train a shared model while
each keeps its own data in-house. We asked which federated method to use, and
whether the approach works outside a simulation. Treating four separate
diabetes studies as collaborating sites, we compared the main federated
strategies and tested each on three further studies that contributed no training
data.
We found that federation was about as accurate as training on a single cohort
for patients from the participating sites, and more reliable on the unseen
cohort: a model trained on one cohort alone sometimes failed badly on new
patients, whereas the federated models did not, on the three held-out cohorts we
tested. The accuracy differences are
small, so the main benefit is lower risk of a bad surprise on a new population
rather than a large accuracy gain. We also built a minimal system and ran the
training for real across two separate computers over the internet. Our results
give a small group of clinics a practical, privacy-preserving starting point for
building glucose predictors that are less likely to fail on patients they have
never seen --- though clinical metrics beyond error size, and tests on more
diverse populations, are still needed before deployment.

\clearpage
\newgeometry{top=0.85in,left=1in,right=1in,footskip=0.75in}
\linenumbers

% ===========================================================================
\section*{Introduction}

Type~1 diabetes (T1D) affects an estimated 8.4~million people worldwide and its
incidence is rising in many regions~\cite{idf2021,gregory2022t1dindex}. People living with T1D must
continually balance insulin, food, and activity to keep blood glucose within a
narrow target range; sustained excursions above or below that range drive both
the acute dangers of hypoglycaemia and the long-term microvascular and
cardiovascular complications of hyperglycaemia. Continuous glucose monitoring
(CGM), which reports interstitial glucose roughly every five minutes, is now
central to diabetes care; it underpins clinical outcomes such as time in range,
for which international consensus targets have been defined~\cite{battelino2019tir}. The same dense data stream makes
short-horizon glucose \emph{forecasting} clinically valuable: a reliable
prediction of where glucose will be in 30--60 minutes can trigger preemptive
alerts, inform bolus and snack decisions, and close the loop in automated
insulin-delivery systems~\cite{vettoretti2020advanced,bekiari2018ap}.

Short-horizon CGM forecasting has progressed from autoregressive and classical
machine-learning models to deep sequence models. Recurrent architectures with
predictive variance estimation established strong baselines on benchmark
cohorts~\cite{martinsson2020}, and standardised datasets and challenges shaped
much of the early evaluation~\cite{marling2020ohio}. The
transformer~\cite{vaswani2017attention} and its long-horizon time-series
variants~\cite{zhou2021informer,wu2021autoformer} have since been adapted to
glucose, with personalised, uncertainty-aware transformer forecasters reporting
state-of-the-art accuracy~\cite{sergazinov2023gluformer} and curated CGM
benchmarks standardising their evaluation~\cite{sergazinov2024glucobench}.
These models are accurate when trained and evaluated on the same cohort, but how
well they transfer across cohorts is an open question~\cite{ghimire2024generalize}:
glucose dynamics, demographics, and devices differ from site to site, and a model
can degrade on populations it did not see in training~\cite{farahmand2025glumind}.
A model that has only seen one cohort therefore has little reason to generalise to
the next; robust generalisation instead calls for training over large, diverse
populations spanning many sites.

Large, diverse training data is exactly what is hardest to assemble. CGM traces
are sensitive personal health records, fragmented across hospitals, clinical
trials, device manufacturers, and national registries, and governed by
data-protection regimes such as the GDPR and HIPAA that restrict moving raw
records between institutions, let alone across borders. The data that improve
generalisation and the data that exist are thus mismatched: large, heterogeneous
corpora are what is needed, but what is available is many small, governed silos.
Federated learning (FL) addresses this tension~\cite{mcmahan2017fedavg}: sites
train on their own data and exchange only model updates, so a shared model is
learned without any raw record leaving its institution. On this basis FL has been
widely adopted for collaborative medical machine
learning~\cite{rieke2020future,sheller2020federated,kaissis2020secure}. Its
central difficulty is statistical heterogeneity: plain federated averaging
degrades when client distributions diverge, motivating proximal
regularisation~\cite{li2020fedprox} and control-variate variance
reduction~\cite{karimireddy2020scaffold}. Most of this work, however, is evaluated
in simulation, where ``clients'' are partitions of one dataset on a single
machine.

Even with federation, real CGM cohorts differ in device generation, demographics,
treatment protocols, and glycaemic regime, and these differences widen across
countries. A federated model must therefore cope with substantial domain shift: it
should generalise to patients, and even whole cohorts, it never trained on (domain
generalisation), and/or specialise to each site (personalisation). Domain
generalisation seeks models that transfer to unseen
distributions~\cite{zhou2022dgsurvey}, and meta-learning for domain generalisation
(MLDG) simulates train/test domain shift within each update so the model learns to
generalise rather than memorise~\cite{li2018mldg}. Bringing these ideas into
federation is an active area surveyed by Li et al.~\cite{li2025fedDGsurvey},
spanning several families. Meta-learning approaches adapt the model-agnostic
meta-learning objective to federation, for communication-efficient
convergence~\cite{chen2018fedmeta} or for models that personalise to a new client
in a few gradient steps~\cite{fallah2020perfedavg}. Personalisation approaches keep
a per-client model: APFL learns an adaptive convex mixture of the global and a
local model~\cite{deng2020apfl}, while Ditto regularises each client's personal
model toward the global one with a tunable proximal strength~\cite{li2021ditto}.
Other approaches instead align client representations, via model-contrastive
objectives~\cite{li2021moon}, federated contrastive learning for multi-centre
medical imaging~\cite{fedcl2023}, or adversarial domain
alignment~\cite{flda2023,fedadv2022}. Our benchmark covers the first two
families---meta-learning (MLDG) and personalisation (APFL, Ditto)---alongside
standard federated averaging (FedAvg, FedProx) and single-cohort local training,
and leaves representation-alignment methods to future work.

\textbf{Yet these methods are rarely compared head-to-head for glucose
forecasting.} They are seldom benchmarked against one another on a common task, and
evaluations rarely separate the two regimes that matter clinically---accuracy on
patients from the training cohorts (held-in) versus accuracy on entirely unseen
cohorts (out of distribution, OOD)---even though a method can behave quite
differently in the two. Closing this gap---comparing federated strategies on one
glucose task under a shared protocol, with held-in and OOD performance reported
separately---is the main focus of this paper.

A second, more practical observation is that almost all federated glucose work is
simulated rather than deployed, with ``clients'' as partitions of one dataset on a
single machine. Running a federation for real additionally requires networked,
physically separate sites to exchange updates and to start from a common model;
production FL frameworks can do this but carry substantial dependencies and
assumptions that can be ill-suited to a small clinical consortium that simply wants
a handful of sites, one cohort each, to train together. We therefore also ask a
feasibility question: whether a lightweight cross-site federation for CGM
forecasting can be built and run across real machines with minimal tooling.

This paper makes one primary contribution and one secondary one. \textbf{Our
primary contribution is a controlled domain-generalisation benchmark of federated
strategies for transformer-based glucose forecasting.} Using a common
encoder--decoder transformer and four real adult T1D cohorts as federated
clients, with a held-out cohort as an OOD test site, we compare single-cohort
local training, FedAvg, FedProx, MLDG, and two personalised-FL methods (APFL and
Ditto) under a multi-seed protocol, reporting held-in and OOD accuracy separately
at the 30-minute horizon --- a horizon commonly used for preemptive glucose
alerts. All cohorts are adult T1D populations, so our findings should not be
assumed to extend to paediatric or type-2 diabetes.

\textbf{As a secondary contribution, we show that such a federation can be run for
real, not only simulated.} We build, from scratch, a lightweight cross-site FL
system for CGM forecasting: a stateless aggregation server that holds the global
model but no data, and clients that each wrap one institution's cohort,
communicating over HTTP using only the Python standard library. The server
broadcasts the initial weights so every client starts from an identical model, and
the same protocol carries FedAvg, FedProx, and MLDG training. We validate it
end-to-end across two physically separate machines connected over the public
internet as a feasibility prototype, and describe how the design could extend to
a geographically distributed, multi-country deployment. We do not claim a
production-grade or multi-country deployment here: establishing that would
require secure aggregation, encryption, and more than two independent sites,
which we leave to future work.

Taken together, the benchmark shows how to evaluate federated glucose
forecasting honestly and what such an evaluation reveals, while the prototype
shows that a federation of this kind can be run across physically separate
machines, not only simulated.

% ===========================================================================
\section*{Materials and methods}

\subsection*{Datasets and federation setup}

% NOTE (author): patient counts here are post-filtering subsets and differ from
% the enrolled totals in the source papers (HUPA-UCM 22 vs 25; T1D-UOM 14 vs 17,
% of 21 enrolled). State the exact inclusion/exclusion criterion at submission.
% See reviews/UNSURE.md item 4.
We use seven publicly available type-1 diabetes (T1D) continuous glucose
monitoring (CGM) datasets. Four serve as the federated clients --- HUPA-UCM
\cite{hidalgo2024hupaucm}, ABC4D \cite{abc4d_nct}, ARISES \cite{zhu2022arises},
and T1D-UOM \cite{alsuhaymi2025t1duom} --- with \emph{one dataset per client}. Each dataset
originates from a distinct study, population, and device, so a per-dataset
partition makes every client a genuinely different domain rather than an
arbitrary shard of a pooled corpus; this is what lets the federation function as
a test of cross-domain generalisation. The remaining three datasets ---
ReplaceBG \cite{aleppo2017replacebg}, BrisT1D, and Flair --- are never seen
during training and are used only as held-out, out-of-distribution (OOD) test
cohorts, so that transfer is measured on three independent populations that
contributed no data to training rather than on a single held-out cohort. Within
each client, we partition patients
into disjoint train, validation, and test groups, splitting strictly by
patient ID so that no patient contributes data to more than one group. Every
window in the validation and test sets therefore comes from a patient never
seen during training, and reported performance measures generalisation to
unseen patients rather than to unseen windows of patients already observed.

Two of the four client cohorts, HUPA-UCM and T1D-UOM, together with the three
held-out OOD cohorts (ReplaceBG, BrisT1D, Flair), are obtained through MetaboNet
\cite{wolff2026metabonet}, a recently released resource that consolidates many
public T1D datasets into a single harmonised schema with a common 5-minute grid
and unified CGM, insulin, and carbohydrate fields. Drawing these cohorts from
MetaboNet ensures their variables are already aligned to a shared format; the
remaining two client cohorts, ABC4D and ARISES, are converted into the same
schema from their original releases. We do not adopt MetaboNet's own train/test split,
which allows the same patient to appear in both partitions; instead we pool the
data and re-split strictly by patient.

Table~\ref{tab:datasets} summarises each cohort. Mean and standard deviation of
CGM are reported in mg/dL, and time-in-range (TIR) is the fraction of CGM
readings in the clinical 70--180 mg/dL band.

\begin{table}[!ht]
\centering
\captionsetup{justification=centering,singlelinecheck=on}
\caption{\bf Federated client cohorts.}
\label{tab:datasets}
\begin{tabular}{lrrrrr}
\thickhline
dataset & patients & total windows & mean CGM & sd CGM & TIR (\%) \\
        &          &               & (mg/dL)  & (mg/dL) & 70--180  \\
\hline
HUPA-UCM & 22 & 291,728 & 140.9 & 56.6 & 72.2 \\
ABC4D    & 25 & 562,043 & 155.1 & 63.4 & 65.5 \\
ARISES   & 12 & 126,248 & 160.5 & 63.7 & 63.7 \\
T1D-UOM  & 14 & 300,796 & 151.6 & 59.9 & 70.5 \\
\thickhline
\end{tabular}
\begin{center}
{\small The three OOD test cohorts (ReplaceBG, BrisT1D, Flair) are excluded from
training and reported separately (Table~\ref{tab:ood}). ARISES is both the
highest-glucose and the smallest client cohort. Patient counts are the usable
subsets after quality filtering and are smaller than the totals enrolled in the
original dataset releases.}
\end{center}
\end{table}

\subsection*{Preprocessing and windowing}

Every cohort is resampled to a uniform 5-minute grid. We form supervised
examples with a sliding window of 72 samples (6\,h) of CGM context to forecast
the next 12 samples (60\,min), advancing the window one sample at a time
(stride~1). Windows that would bridge a sensor gap longer than 6 samples
(30\,min) are discarded, so every example is contiguous. CGM is normalised with
a single \emph{global} mean and standard deviation ($154.04 \pm 61.00$ mg/dL)
computed from the four training cohorts only --- the held-out OOD cohort does not
contribute to it, so no OOD information enters preprocessing --- and shared
across all clients rather than using per-client statistics. A common
normalisation keeps the global model's input scale identical at every site,
which is what allows weights aggregated across clients to be applied directly to
the OOD cohorts; it is shipped with the model as a fixed constant. Alongside CGM,
each example carries a time-of-day mark encoded
as $(\sin 2\pi\,\mathrm{frac},\ \cos 2\pi\,\mathrm{frac})$, where $\mathrm{frac}$
is the fraction of the day, to expose circadian structure.

\subsection*{Forecasting model}

The forecaster is a sequence-to-sequence (encoder--decoder) Transformer that maps
the 72-step CGM context to the 12-step horizon. It uses
$d_{\text{model}}=128$, 8 attention heads, 6 encoder layers, 2 decoder layers,
feed-forward width $d_{\text{ff}}=2048$, ReLU activations, and dropout~$0.05$
($\approx 4.9$\,M parameters). Decoding is single-shot rather than
autoregressive: the decoder is fed a learned start token while the encoder
carries the full history, and the 12-step horizon is produced in one forward
pass. The decoder predicts a single CGM channel ($c_{\text{out}}=1$) and the
model is trained with mean-squared-error (MSE) loss. The same architecture is
used for every client and every federated strategy, so all comparisons differ
only in how the model is trained, not in capacity. The model emits a point
forecast only; it produces no predictive-uncertainty estimate (see Limitations).

\paragraph{Bolus input variant}
As a secondary analysis, we test whether adding an exogenous input improves
forecasting. We train a two-channel variant ($\texttt{enc\_in}=2$) in which a
second encoder channel carries $\log(1+\mathrm{bolus}_{\mathrm{IU}})$, the
log-transformed bolus insulin dose; bolus events are sparse (roughly 1\% of
timesteps are non-zero). The decoder remains CGM-only. We rerun the full
multi-seed sweep in this two-channel configuration and report it as a delta
against the CGM-only model (Results); because the effect on the global models
proves small relative to seed-to-seed variation, we treat it as a sensitivity
check rather than a main result.

\subsection*{Federated learning system}

Training is coordinated by a from-scratch, networked federated learning system
in which one institution runs a server and each cohort runs a client. The server
holds the global model weights but never sees any data: at round one it mints a
single seeded set of initial weights and broadcasts them, so every client begins
from byte-identical parameters regardless of hardware. Each round, clients fetch
the current global weights, train locally on their own cohort, and return updated
weights; the server aggregates them by FedAvg and advances to the next round.
Clients communicate with the server over plain HTTP through a small task
state-machine (train\,$\rightarrow$\,validate\,$\rightarrow$\,test), so the same
protocol carries FedAvg, FedProx, and MLDG without modification --- the choice of
strategy is entirely client-side, and the server always performs the same FedAvg
aggregation. We validate the system end-to-end across two physically separate
machines communicating over the public internet as a feasibility prototype, and
describe how the design could generalise to a geographically distributed,
multi-country deployment. The prototype uses plain HTTP without encryption,
secure aggregation, or differential privacy; it demonstrates the training
workflow, not a production-grade secure channel (see Limitations).

\subsection*{Federated strategies and baselines}

\textbf{Local (baseline).} Our primary comparator is local training: a separate
model trained on a single cohort's data only, with no federation. The central
question of the benchmark is whether federation improves on this baseline ---
both when evaluated on the same cohorts that participated (held-in) and on the
unseen OOD cohort. For OOD we report the mean over the four single-cohort models.

\textbf{FedAvg / FedProx.} Standard federated averaging, and FedProx, which adds
a proximal term $(\mu/2)\lVert w - w_{\text{global}}\rVert^2$ to each client's
objective to limit client drift. We run FedProx at a single proximal strength
($\mu=0.05$) and did not sweep it, so its ranking relative to FedAvg should be
read as indicative rather than as a tuned comparison.

\textbf{MLDG.} Meta-learning for domain generalisation, adapted to the federated
setting: at each client step the client's patient groups are split into
meta-train and meta-test partitions, an inner adaptation step is taken on
meta-train, and the model is updated by the gradient of the meta-test loss
evaluated after that step. Because the split is over patients \emph{within} a
client, MLDG here simulates cross-\emph{patient} shift inside a cohort, not the
cross-\emph{cohort} shift measured by the OOD test. MLDG and FedProx alter only
client-side training; the server still FedAvg-aggregates. Concretely, each client
step takes one inner gradient step on the meta-train patients, evaluates the loss
on the held-out meta-test patients after that step, and updates the model by the
gradient of that meta-test loss; the exact inner learning rate and the handling
of the optimiser state across the inner and outer steps are set in configuration
and will be released with the code (see Data and code availability).

\textbf{Personalised FL (PFL).} APFL learns a per-client mixture
$\alpha v_i + (1-\alpha)w$ of a personal model $v_i$ and the global model $w$
with an adaptively updated $\alpha$ (initial $0.25$); we also include a decoupled
APFL variant. Ditto trains a personal model pulled toward the global one by a
proximal term ($\mu\in\{0.01,0.1,1.0\}$). These methods produce one personalised
model per client and \emph{no} single global model, so --- like the local
baseline's per-cohort models --- they are evaluated held-in only and cannot be
applied to the OOD cohorts. We flag that the APFL mixing weight may not have moved
far from its initial value in our runs; its numbers should therefore be read as a
lower bound on what a fully tuned APFL could achieve.

\subsection*{Evaluation protocol and metrics}

We evaluate every strategy in two regimes. \emph{Held-in} performance is measured
on each participating cohort's own held-out patients, isolating predictive
accuracy at the contributing sites. \emph{OOD} performance is measured on three
never-seen cohorts (ReplaceBG, BrisT1D, Flair), which contributed no training
data, isolating cross-domain transfer; only strategies that yield a single
deployable global model (Local-mean, FedAvg, FedProx, MLDG) receive an OOD score.
We report root-mean-squared error (RMSE) in mg/dL at the 30-minute horizon; the
model forecasts to 60 minutes, but we defer the longer-horizon analysis to future
work. All runs use early stopping on average validation MSE (patience~5 over a
maximum of 25 communication rounds), and every configuration is repeated over 5
random seeds (42, 43, 44, 46, 47; seed 45 was not run). We report the mean and
standard deviation across seeds; the global and single-cohort runs are complete
over all five seeds, while a few personalised-FL held-in cells use four seeds
where a run is missing. To test whether the small differences between federated
strategies are meaningful, we compare per-seed RMSE with a paired $t$-test across
the common seeds. We emphasise that these standard deviations quantify
\emph{optimisation stochasticity} for one fixed patient split; because we use a
single by-patient split, they understate the patient-sampling uncertainty that
would matter at deployment, and should not be read as precision on new patients.
Alongside RMSE we report time-in-range agreement and forecast lag as
clinically oriented metrics; an error-grid analysis, MARD, and
hypo-/hyperglycaemia event detection are not included here and are discussed as
limitations.

\subsection*{Implementation and training details}

All models are trained with the Adam optimiser (learning rate~$10^{-4}$, weight
decay~$10^{-3}$, batch size~256, one local epoch per communication round). The
federated system is implemented in Python using only the standard library for
networking.

\subsection*{Data and code availability, and ethics}

All seven cohorts are publicly available, de-identified secondary datasets, used
here under their respective terms of access; no new human-subjects data were
collected for this study. HUPA-UCM \cite{hidalgo2024hupaucm}, T1D-UOM
\cite{alsuhaymi2025t1duom}, and the three OOD cohorts ReplaceBG
\cite{aleppo2017replacebg}, BrisT1D, and Flair are obtained via MetaboNet
\cite{wolff2026metabonet}; ABC4D and ARISES are used from their original releases. The code for the federated system and the analysis, together
with the exact configuration and seeds behind each table, will be released in a
public repository upon publication.
% NOTE (author): insert the repository DOI/URL here at submission, and add each
% dataset's access route/licence/accession (Zenodo/NCT) plus an IRB/consent
% statement for the secondary use. Also add the FL-architecture figure. See
% reviews/UNSURE.md items 5, 20.

\section*{Results}

Table~\ref{tab:main} reports held-in RMSE on each client's unseen test patients
and OOD RMSE on ReplaceBG, at the 30-minute horizon, for the CGM-only model.

\begin{table}[!ht]
\centering
\captionsetup{justification=centering,singlelinecheck=on}
\caption{\bf CGM-only RMSE (mg/dL) at the 30-minute horizon.}
\label{tab:main}
\begin{tabular}{lrrrrrr}
\thickhline
 & \multicolumn{5}{c}{held-in (per client)} & OOD \\
\cline{2-6}
strategy & HUPA-UCM & ABC4D & ARISES & T1D-UOM & avg & ReplaceBG \\
\hline
Local (per cohort) & 18.81 & 19.95 & 22.04 & 19.80 & 20.15 & 23.54 \\
\hline
FedAvg  & 18.64 & 19.80 & 22.38 & 19.53 & 20.09 & 21.43 \\
FedProx & 18.89 & 19.94 & 22.60 & 19.67 & 20.28 & 21.57 \\
MLDG    & \textbf{18.46} & \textbf{19.68} & 22.30 & 19.53 & \textbf{19.99} & \textbf{21.34} \\
\hline
APFL          & 19.11 & 19.88 & 22.53 & 19.61 & 20.28 & --- \\
APFL-decoupled & 19.37 & 19.90 & 22.47 & 19.73 & 20.37 & --- \\
Ditto ($\mu{=}0.01$) & 18.90 & 19.90 & \textbf{21.98} & 19.64 & 20.11 & --- \\
Ditto ($\mu{=}0.1$)  & 18.90 & 19.93 & 22.28 & 19.66 & 20.19 & --- \\
Ditto ($\mu{=}1.0$)  & 18.74 & 19.77 & 22.35 & \textbf{19.50} & 20.09 & --- \\
\thickhline
\end{tabular}
\begin{center}
{\small Mean over 5 random seeds (42, 43, 44, 46, 47) for the global and Local
rows; the personalised (APFL, Ditto) held-in rows use 4 seeds where a run is
missing. Standard deviations are across seeds and are small ($\approx 0.1$--$0.3$
mg/dL held-in and OOD), except the Local OOD entry, which is noisy across seeds
(sd $\approx 1.8$ mg/dL). \emph{avg} is the \emph{unweighted} mean of the four
held-in client means (cohorts differ in size, so it is not a window-weighted
average). Local held-in is each cohort's own single-cohort model; Local OOD
($23.54$) is the mean of the four single-cohort models on ReplaceBG.
Table~\ref{tab:ood} extends the OOD comparison to two further held-out cohorts.
Personalised strategies (APFL, Ditto) yield no single global model and so receive
no OOD score. Bold marks the numerically lowest value per column; differences
among the federated strategies are within seed variation (see text).}
\end{center}
\end{table}

\paragraph{Federation matches local accuracy held-in and transfers more reliably
out-of-distribution, on all three unseen cohorts.} On the contributing cohorts,
the global FedAvg and MLDG models are on par with cohort-specific local training
(held-in average $20.09$ and $19.99$ vs.\ $20.15$ mg/dL): pooling four
heterogeneous cohorts through federation does not cost in-domain accuracy despite
each client seeing only the shared global weights. The difference emerges
out-of-distribution, and it is about reliability more than average accuracy. We
test transfer on three independent held-out cohorts --- ReplaceBG, BrisT1D, and
Flair (Table~\ref{tab:ood}) --- and the pattern is the same on all three: every
global strategy beats the mean single-cohort model, and the mean single-cohort
model is dragged up by one failure. On ReplaceBG the federated models scored
$21.34$--$21.57$ mg/dL; three of the four single-cohort models transferred to
within roughly half a mg/dL of them, while the ARISES-only model collapsed
($28.51 \pm 6.85$ mg/dL). The same single model collapsed on BrisT1D
($32.39 \pm 5.84$) and Flair ($30.99 \pm 6.23$), with a seed-to-seed standard
deviation an order of magnitude larger than any federated model's. The honest
reading is therefore not that federation is a couple of mg/dL more accurate ---
a gap that is in any case small relative to CGM sensor error
($\approx 8$--$15$ mg/dL) --- but that a single-cohort model can transfer
catastrophically and one cannot know in advance which cohort's model will fail,
whereas across three unseen cohorts no federated model did so. Federation's value
here is reduced tail risk, not a large mean gain.

\begin{table}[!ht]
\centering
\captionsetup{justification=centering,singlelinecheck=on}
\caption{\bf OOD RMSE (mg/dL) at 30 min on three held-out cohorts (CGM-only).}
\label{tab:ood}
\begin{tabular}{lrrr}
\thickhline
strategy & ReplaceBG & BrisT1D & Flair \\
\hline
FedAvg  & $21.43 \pm 0.08$ & $26.37 \pm 0.07$ & $25.07 \pm 0.09$ \\
FedProx & $21.57 \pm 0.06$ & $26.51 \pm 0.13$ & $25.22 \pm 0.11$ \\
MLDG    & $\mathbf{21.34 \pm 0.17}$ & $\mathbf{26.31 \pm 0.21}$ & $\mathbf{25.05 \pm 0.15}$ \\
\hline
Local (mean of 4 singles) & ${\approx}23.5$ & ${\approx}28.1$ & ${\approx}26.8$ \\
\;\emph{of which} single-ARISES & $28.51 \pm 6.85$ & $32.39 \pm 5.84$ & $30.99 \pm 6.23$ \\
\thickhline
\end{tabular}
\begin{center}
{\small Mean $\pm$ sd over 5 seeds. On every cohort all three federated
strategies beat the mean single-cohort model, and MLDG is numerically lowest;
the single-cohort model trained on ARISES (the smallest cohort) collapses on all
three, with a seed-to-seed sd $30$--$40\times$ that of any federated model. The
Local mean is the unweighted mean of the four single-cohort models (its exact
value shifts by $<0.3$ mg/dL with how missing single-model seeds are averaged;
the single-ARISES term dominates its variance regardless).}
\end{center}
\end{table}

\paragraph{MLDG is a small, consistent, but not statistically significant
improvement.} Among the strategies that produce a deployable global model, MLDG
has the numerically lowest error in \emph{every} held-in and OOD comparison, and
it beats FedAvg on 4 of the 5 seeds in each regime. The effect is directionally
consistent but small: the held-in-average gap is $0.09$ mg/dL and the OOD gaps
are $0.02$--$0.09$ mg/dL, and a paired test across seeds does not reach
significance ($t \approx -1.3$, $n=5$ for the held-in average). We therefore
describe MLDG as a small, consistent advantage rather than a definitive one, and
we do not claim a statistically significant ranking among the three global
strategies. FedProx was numerically \emph{weakest} of the three (highest RMSE
held-in and on all three OOD cohorts), but it was run at a single untuned $\mu$,
so this is not a tuned ranking. The personalised methods do not change the
held-in picture --- their averages cluster around FedAvg's ($20.1$--$20.4$) ---
and by construction cannot be deployed on the OOD cohorts.

\paragraph{A federated model is a competitive starting point for a new site.} We
next asked whether the federated model helps a \emph{new} site that does
contribute data, by finetuning the global model on ReplaceBG's own training
patients and comparing against a model trained on ReplaceBG from scratch.
Finetuning any of the three global models reached a slightly lower RMSE on
ReplaceBG's held-out patients (FedProx$\rightarrow$finetune $20.98 \pm 0.06$,
FedAvg$\rightarrow$finetune $21.04 \pm 0.04$, MLDG$\rightarrow$finetune
$21.04 \pm 0.07$ mg/dL) than training on ReplaceBG alone
($21.38 \pm 0.09$ mg/dL). The differences are small ($\approx 0.3$--$0.4$ mg/dL)
and we report no significance test, so we read this as suggestive: the federated
model is at least a competitive initialisation, and training on the new site's
own data did not beat it. We note two caveats that bound this result: ReplaceBG
is a large cohort ($166$ training patients), so this does not test the
data-poor small-clinic regime the deployment argument targets, and the
from-scratch comparator used the same training budget, which we did not
separately tune.

\paragraph{Beyond RMSE: time-in-range agreement and prediction lag.} As a first
step toward clinically oriented metrics, we also report, for the global models,
how well predicted time-in-range (TIR, 70--180 mg/dL) matches actual TIR and how
much the 30-minute forecast lags the true trajectory. Predicted and actual TIR
agree to within a few percentage points --- for MLDG, $3.6$ points held-in
(pooled over cohorts) and $1.7$--$2.7$ points on the three OOD cohorts --- and
FedAvg is similar, so the federated models preserve population-level glycaemic
statistics on unseen cohorts. The cross-correlation lag of the forecast is
$\approx 20$ minutes held-in and $\approx 20$--$22$ minutes OOD, i.e.\ a
non-trivial fraction of the 30-minute horizon is effectively recent-value
persistence --- a known behaviour of CGM forecasters and a reason a point-RMSE
improvement need not translate into earlier warnings. We did \emph{not} compute a
Clarke/Parkes error grid or event-level hypo-/hyperglycaemia detection, which
require the per-window predictions and are the clearest next addition for a
clinical venue (see Limitations).

\paragraph{Adding bolus insulin does not help the federated models, but it damps
the single-cohort collapse.} We reran the full sweep with a second encoder
channel carrying log-bolus insulin. For the global models the effect is
essentially null: held-in and OOD RMSE move by less than $\approx 0.1$ mg/dL in
either direction (within seed variation), and MLDG's OOD ReplaceBG error is
unchanged at $21.34$ mg/dL. The one place bolus makes a consistent difference is
the failure mode itself: the single-ARISES model's OOD collapse is milder and far
less variable with bolus than without it, on all three held-out cohorts
(ReplaceBG $26.86 \pm 3.45$ vs $28.51 \pm 6.85$; BrisT1D $31.95 \pm 3.42$ vs
$32.39 \pm 5.84$; Flair $29.81 \pm 3.09$ vs $30.99 \pm 6.23$ --- the standard
deviation roughly halves in each case). The natural reading is that an exogenous
input acts as a mild regulariser against the worst single-cohort overfitting ---
but this is precisely the failure that federation already prevents, so bolus adds
little once a model is federated. We therefore treat bolus as a sensitivity
check, not a main result.

\section*{Discussion}

\paragraph{Principal findings.} Across four real adult T1D cohorts federated as
distinct domains, with three never-seen cohorts as out-of-distribution (OOD)
tests, federation matched single-cohort local training in-domain and transferred
more reliably out of it. Held-in, the global FedAvg and MLDG models were on par
with cohort-specific local training (average RMSE $20.09$ and $19.99$ vs.\
$20.15$ mg/dL at the 30-minute horizon); OOD, on all three held-out cohorts the
federated models beat the mean single-cohort model and avoided the catastrophic
transfer failure that the smallest-cohort model showed on every one of them ---
though most single-cohort models transferred comparably and the mean gain is
small relative to sensor error. Among the federated strategies MLDG was
numerically best in every comparison but by margins within seed-to-seed
variation. Beyond the benchmark, we showed that a federation of this kind runs
across physically separate machines: the same protocol trains FedAvg, FedProx,
and MLDG over the public internet, with no raw data leaving any site --- a working
artifact, not a simulation.

\paragraph{Federation buys more reliable transfer, not only an accuracy gain.}
The clearest message of the benchmark is not that federation is more accurate
in-domain --- the held-in gap over local training is small --- but that it
lowers the risk of a catastrophic transfer failure that single-cohort training
carries. A model trained on one cohort can transfer badly: the ARISES-only
model, which held-in was on par with federation on its own patients ($22.04$ vs
MLDG $22.30$ mg/dL, this being the hardest cohort for every method), collapsed on
\emph{every} held-out cohort ($28.51 \pm 6.85$ on ReplaceBG, $32.39 \pm 5.84$ on
BrisT1D, $30.99 \pm 6.23$ on Flair) with a seed-to-seed variance far larger than
any federated model's, even though the other three single-cohort models
transferred within about half a mg/dL of the federated models. The difficulty is
that one cannot know in advance which site's model will be the one that fails.
Across three independent held-out cohorts, every federated model avoided this
failure --- a pattern replicated rather than observed once, which materially
strengthens the claim over a single-cohort result. We still stop short of calling
it a guarantee (three adult cohorts is not the full space of distribution shift,
and a leave-one-cohort-out design would test it further; see Limitations), but
for a clinical setting in which distribution shift is inevitable, reducing the
chance of a silent, cohort-specific collapse is a practical argument for
federation independent of any in-domain accuracy difference.

\paragraph{Generalisation versus personalisation, and how to have both.} Our
results expose a genuine tension. The personalised strategies (Ditto, APFL) and
dedicated single-cohort models were the \emph{only} approaches that beat the
federated models on the hardest and smallest in-domain cohort: on ARISES, Ditto
($\mu{=}0.01$, $21.98$ mg/dL) and the dedicated ARISES model ($22.04$ mg/dL) edged
MLDG ($22.30$). Yet these same methods produce no single global model and so
cannot serve the OOD cohorts at all --- exactly the regime in which the global
models are most valuable. The two objectives pull in opposite directions:
specialising to a site helps that site but forfeits transfer; generalising across
sites transfers but leaves per-site accuracy on the table. The way to obtain both
is to treat the generalised global model as a foundation and adapt it locally,
rather than choosing one regime outright.

\paragraph{Finetuning a federated model for a new site.} The finetuning
experiment points, tentatively, toward a ``foundation-then-adapt'' recipe.
Finetuning the global federated model on ReplaceBG's own patients reached
$20.98$--$21.04$ mg/dL, slightly better than either the zero-shot federated model
($21.34$ mg/dL) or a model trained on ReplaceBG from scratch ($21.38$ mg/dL). The
differences are small ($\approx 0.3$--$0.4$ mg/dL) and untested for significance,
so we draw only qualitative conclusions. First, training on the site's own data
in isolation did not beat the finetuned federated model, suggesting the shared
model contributes usefully. Second, the three global strategies converge once
finetuned: FedProx, numerically \emph{weakest} zero-shot here, is numerically
\emph{best} after finetuning ($20.98$ mg/dL) --- a reversal consistent with the
choice of federated strategy mattering most for zero-shot transfer and mattering
little once local data are available. A
plausible practical protocol is therefore two-stage --- federate to obtain a
shared model, then finetune per site --- but establishing that a finetuned model
reliably beats local-only training, especially for a data-poor new site, requires
a data-efficiency study we have not yet run (this experiment used ReplaceBG's
large 166-patient training set).

\paragraph{On meta-learning.} MLDG was numerically best in every held-in and OOD
comparison and beat FedAvg on four of five seeds --- a strikingly consistent
pattern --- yet the per-comparison margins ($\le 0.1$ mg/dL) are within seed
noise and a paired test at five seeds does not reach significance, so we report a
small, consistent edge rather than a demonstrated meta-learning advantage. A
mechanistic caveat tempers the interpretation further: our MLDG splits patients
into meta-train and meta-test \emph{within} each client, so it simulates
cross-\emph{patient} shift inside a cohort, not the cross-\emph{cohort} shift the
OOD test measures. That it nonetheless comes out consistently ahead OOD is
suggestive but unexplained; a federated-DG formulation that meta-splits
\emph{across} clients would target the evaluated shift more directly and, with
more seeds to power a test, is the natural next step. The dominant effect in our
results is federation itself; the choice of federated strategy is second-order.

\paragraph{Implications for a clinical consortium.} These findings give a small
group of institutions a reason to consider federating. A participating centre
loses little in-domain, and gains a model less likely to fail on a population it
did not train on. Because no raw record leaves any institution, federation
removes the raw-data-pooling barrier that data-protection regimes such as GDPR
and HIPAA impose --- though we stress that keeping raw data in-house is a
\emph{necessary but not sufficient} condition for compliance: a real deployment
would still need encryption in transit, secure aggregation or differential
privacy against inversion of the exchanged updates, access control, and an
appropriate governance and impact assessment, none of which our prototype
provides. Subject to those safeguards, the tooling itself is light enough to run
without dedicated federated-learning infrastructure. The benefit should also
compound as more cohorts join and the shared model becomes more representative.

\paragraph{Strengths.} Three features distinguish this study from prior federated
glucose work. First, it compares local training, FedAvg, FedProx, MLDG, and two
personalised-FL families head-to-head on one task under a shared protocol, rather
than reporting a single method in isolation. Second, it separates held-in from
OOD performance, a distinction that prior evaluations rarely make explicit but
that proves informative: the strategies are nearly indistinguishable in-domain,
and OOD the meaningful split is between single-cohort training (which can fail)
and federation (which did not, here), rather than between federated strategies.
Third, the cohorts are partitioned strictly by patient and normalised with a
single shared statistic, and the federation is run across real, physically
separate machines rather than simulated as shards of one dataset on one host.

\paragraph{Positioning against prior results on these cohorts.} Our per-cohort
accuracy is competitive with previously published results on the same datasets,
which is notable because we evaluate under a strictly harder protocol. The
distinction is the train/test split. Almost all prior work splits each
patient's series chronologically --- the first portion for training, the last
portion for testing (a \emph{by-time} split) --- so every test window comes from
a patient the model has already seen; the model need only extrapolate forward in
time for a familiar individual. We instead split \emph{by patient}: test patients
are disjoint from training patients, so every reported number measures
generalisation to people the model has never seen. By-patient evaluation is the
harder task and the one that matches deployment, where a predictor meets new
patients, but it typically yields higher error than a by-time split on the same
cohort. Table~\ref{tab:prior} collects the closest comparable external results.
On HUPA-UCM and T1D-UOM, the strongest real-data RMSE@30 baselines are from
GlucoFM-Bench \cite{lu2026glucofmbench} (a benchmark of time-series foundation
models, univariate CGM, by-time split): HUPA-UCM 16.71 and T1D-UOM 20.81 mg/dL.
On ARISES and ABC4D they are from the personalised evidential meta-learning
forecaster of Zhu et al.\ \cite{zhu2023fcnn} (CGM plus carbohydrate and insulin,
by-time split): 20.23 and 20.25 mg/dL. Our by-patient figures for a global model
(MLDG, numerically the best of our federated strategies; HUPA-UCM 18.46, ABC4D
19.68, ARISES 22.30, T1D-UOM 19.53 mg/dL) sit in the same 17--22 mg/dL range
despite the harder split and, on
HUPA-UCM/T1D-UOM, despite using CGM alone against models that also consume
insulin and carbohydrate inputs. We therefore do not claim to beat these numbers
--- a by-patient result is not expected to --- but to match them under an
evaluation that better reflects deployment. To our knowledge no prior paper
reports all four of these cohorts, separates held-in from out-of-distribution
performance, or evaluates across-patient generalisation as we do; the one prior
work that also splits by patient (AEGIS \cite{pansheriya2026aegis}, on T1D-UOM
via leave-one-patient-out) forecasts only to 15 minutes and so is not comparable
at our 30-minute horizon.

\begin{table}[!ht]
\centering
\captionsetup{justification=centering,singlelinecheck=on}
\caption{\bf Published RMSE@30 on our cohorts, with split protocol.}
\label{tab:prior}
\begin{tabular}{lllll}
\thickhline
study (best model) & cohort & RMSE@30 (mg/dL) & split & sd over \\
\hline
GlucoFM-Bench \cite{lu2026glucofmbench} (LSTM) & HUPA-UCM & 16.71 (4.82) & by-time & patients \\
GlucoFM-Bench \cite{lu2026glucofmbench} (LSTM) & T1D-UOM  & 20.81 (4.26) & by-time & patients \\
Zhu et al.\ \cite{zhu2023fcnn} (FCNN)          & ARISES   & 20.23 (3.38) & by-time & patients \\
Zhu et al.\ \cite{zhu2023fcnn} (FCNN)          & ABC4D    & 20.25 (2.60) & by-time & patients \\
\hline
Ours (MLDG, global)            & HUPA-UCM & 18.46 (0.15) & \emph{by-patient} & seeds \\
Ours (MLDG, global)            & ABC4D    & 19.68 & \emph{by-patient} & seeds \\
Ours (MLDG, global)            & ARISES   & 22.30 & \emph{by-patient} & seeds \\
Ours (MLDG, global)            & T1D-UOM  & 19.53 & \emph{by-patient} & seeds \\
\thickhline
\end{tabular}
\begin{center}
{\small A \emph{by-time} split tests on later windows of patients seen during
training; a \emph{by-patient} split tests on entirely unseen patients and is the
harder, deployment-relevant task. External sd is across patients; ours is across
5 seeds. Numbers are indicative, not head-to-head, because the split (and, for
Zhu et al., the input features) differ. Other papers on these cohorts report
non-comparable metrics or horizons: Tan \& McBeth \cite{tan2026uncertainty}
(MARD only), AEGIS \cite{pansheriya2026aegis} (15-min horizon, mmol/L),
Manchanda et al.\ \cite{manchanda2025lstm} (window-averaged RMSE), and Kolev
et al.\ \cite{kolev2026benchmark} (synthetic data only).}
\end{center}
\end{table}

\paragraph{Limitations.} Several constraints bound these conclusions. \emph{(i)
Evaluation metric.} We report a single aggregate error (RMSE at 30 min). This is
behind the norm of the comparator literature, which reports clinical error grids
(Clarke/Parkes), MARD, or hypo-/hyperglycaemia detection; aggregate RMSE is
dominated by the euglycaemic range and can be near-blind to the safety-critical
hypoglycaemic tail. Our model also emits a point forecast with no calibrated
predictive uncertainty. Whether federation's small advantage survives on
clinically stratified and event-detection metrics is only partly tested (we
report time-in-range agreement and prediction lag, below) and a full error-grid
analysis remains the most important gap for a clinical reading of this work.
\emph{(ii) OOD breadth.} Our transfer claim rests on three held-out cohorts
(ReplaceBG, BrisT1D, Flair) --- a replicated rather than single observation ---
but all are adult cohorts and none stresses paediatric, type-2, or a markedly
different device regime; a leave-one-cohort-out design would probe transfer
further still. \emph{(iii) Statistical power.} Margins among federated strategies
($\le 0.1$ mg/dL) are within seed standard deviation and a paired test at five
seeds does not reach significance, so we do not claim a ranking; MLDG's
consistency (best in every comparison, 4/5 seeds) is suggestive but underpowered.
Seed sd also reflects one fixed patient split and understates patient-sampling
uncertainty. \emph{(iv) Baselines.} We include no persistence or
simple learned baseline in the main table, so we cannot fully attribute accuracy
to the transformer versus the training scheme; FedProx used a single untuned
$\mu$, and APFL may be under-tuned. \emph{(v) Horizon.} We report only 30 min,
though the model forecasts to 60 min. \emph{(vi) Privacy.} Raw data never moves,
but the prototype provides no formal guarantee: it uses plain HTTP without
encryption, secure aggregation, or differential privacy, and the exchanged
updates could in principle be attacked. \emph{(vii) Population.} All seven cohorts
are adult T1D, leaving paediatric and type-2 transfer untested, and we scope the
benchmark to meta-learning and personalisation, leaving representation-alignment
methods to future work. \emph{(viii) No pooled-central control.} We compare
federation only against single-cohort local training, not against a model trained
on all four cohorts pooled on one machine. Federation and single-cohort training
differ in both the algorithm and the amount and diversity of training data, so
strictly our results show that \emph{training across four diverse cohorts} buys
more reliable transfer; whether the federated algorithm matches centralised
pooling (its natural upper bound) is untested. Since federation's purpose is to
obtain that data-diversity benefit \emph{without} pooling raw data, this does not
undercut the practical argument, but it does mean the mechanism should be read as
``more diverse training data, obtained under a privacy constraint'' rather than
anything specific to the federated averaging itself.

\paragraph{Future work.} The next steps follow directly from these limitations.
Most immediately: extend the clinically oriented metrics we begin here (TIR
agreement, forecast lag) to a full Clarke/Parkes error grid, MARD, and event-level
hypo-/hyperglycaemia detection with predictive uncertainty, to test whether
federation's advantage holds where safety matters; add persistence and simple
learned baselines and a centrally pooled control (to separate the federated
algorithm from the diversity of the pooled data); and add more seeds to power the
strategy comparison. A leave-one-cohort-out evaluation would extend transfer
testing beyond the three held-out cohorts reported here. Beyond that: extend to the
60-minute horizon and to event-aware inputs such as insulin and carbohydrate
information; add secure aggregation and differential privacy so the
no-raw-data-movement property is backed by a formal privacy guarantee; and scale
the prototype toward a genuinely multi-site, multi-country deployment with
reported communication, latency, and failure behaviour. Finally, combining the
generalisation of a federated model with per-site finetuning --- ideally trained
jointly, and studied as a function of how much local data a new site has --- and
continual federation as populations and devices drift over time are natural
extensions this framework is positioned to support.


% ===========================================================================
% Compile references.bib with the plos2025.bst style (ships with PLOS template).
\bibliography{references}

\end{document}
